\clearpage
\SourceLeaf{454}%
\pdfbookmark[0]{Gauss an Enke}{gauss-an-enke-1849}
\renewcommand{\BookRunningTitle}{Gauss an Enke}
\begingroup
\begin{center}{\large GAUSS AN ENKE.}\end{center}
\BookRule{.20\textwidth}
\begin{center}Hochzuverehrender Freund!\end{center}

— — Die gütige Mittheilung Ihrer Bemerkungen über die Frequenz der Primzahlen ist mir in mehr als einer Beziehung interessant gewesen. Sie haben mir meine eigenen Beschäftigungen mit demselben Gegenstande in Erinnerung gebracht, deren erste Anfänge in eine sehr entfernte Zeit fallen, ins Jahr 1792 oder 1793, wo ich mir die \textsc{Lambert}’schen Supplemente zu den Logarithmentafeln angeschafft hatte. Es war noch ehe ich mit feineren Untersuchungen aus der höhern Arithmetik mich befasst hatte eines meiner ersten Geschäfte, meine Aufmerksamkeit auf die abnehmende Frequenz der Primzahlen zu richten, zu welchem Zweck ich dieselben in den einzelnen Chiliaden abzählte, und die Resultate auf einem der angehefteten weissen Blätter verzeichnete. Ich erkannte bald, dass unter allen Schwankungen diese Frequenz durchschnittlich nahe dem Logarithmen verkehrt proportional sei, so dass die Anzahl aller Primzahlen unter einer gegebenen Grenze $n$ nahe durch das Integral
\[
\int\frac{\mathrm{d}n}{\log n}
\]
ausgedrückt werde, wenn der hyperbolische Logarithm. verstanden werde. In späterer Zeit, als mir die in \textsc{Vega}’s Tafeln (von 1796) abgedruckte Liste bis 400031 bekannt wurde, dehnte ich meine Abzählung weiter aus, was jenes Verhältniss bestätigte. Eine grosse Freude machte mir 1811 die Erscheinung von \textsc{Chernac}’s
\SourceLeaf{455}%
cribrum, und ich habe (da ich zu einer anhaltenden Abzählung der Reihe nach keine Geduld hatte) sehr oft einzelne unbeschäftigte Viertelstunden verwandt, um bald hie bald dort eine Chiliade abzuzählen; ich liess jedoch zuletzt es ganz liegen, ohne mit der Million ganz fertig zu werden. Erst später benutzte ich \textsc{Goldschmidt}’s Arbeitsamkeit, theils die noch gebliebenen Lücken in der ersten Million auszufüllen, theils nach \textsc{Burck\-hardt}’s Tafeln die Abzählung weiter fortzusetzen So sind (nun schon seit vielen Jahren) die drei ersten Millionen abgezählt, und mit dem Integralwerth verglichen. Ich setze hier nur einen kleinen Extract her:

\begin{center}
\begingroup\fontsize{10}{11.5}\selectfont\setlength{\tabcolsep}{3pt}
\begin{tabular}{r|r|rr|rr}
Unter & \shortstack{gibt es\\Primzahlen} & \shortstack{Integral\\$\displaystyle\int\frac{\mathrm{d}n}{\log n}$} & Abweich. & \shortstack{Ihre\\Formel} & Abweich. \\
500000 & 41556 & $41606{,}4$ & $+50{,}4$ & $41596{,}9$ & $+40{,}9$ \\
1000000 & 78501 & $79627{,}5$ & $+126{,}5$ & $78672{,}7$ & $+171{,}7$ \\
1500000 & 114112 & $114263{,}1$ & $+151{,}1$ & $114374{,}0$ & $+264{,}0$ \\
2000000 & 148883 & $149054{,}8$ & $+171{,}8$ & $149233{,}0$ & $+350{,}0$ \\
2500000 & 183016 & $183245{,}0$ & $+229{,}0$ & $183495{,}1$ & $+479{,}1$ \\
3000000 & 216745 & $216970{,}6$ & $+225{,}6$ & $217308{,}5$ & $+563{,}5$
\end{tabular}\endgroup
\end{center}

Dass \textsc{Legendre} sich auch mit diesem Gegenstande beschäftigt hat, war mir nicht bekannt, auf Veranlassung Ihres Briefes habe ich in seiner \textit{Théorie des Nombres} nachgesehen, und in der zweiten Ausgabe einige darauf bezügliche Seiten gefunden, die ich früher übersehen (oder seitdem vergessen) haben muss. \textsc{Legendre} gebraucht die Formel
\[
\frac{n}{\log n-A}
\]
wo $A$ eine Constante sein soll, für welche er $1{,}08366$ setzt. Nach einer flüchtigen Rechnung finde ich danach in obigen Fällen die Abweichung
\[
\begin{array}{r}
-23{,}3\\
+42{,}2\\
+68{,}1\\
+92{,}8\\
+159{,}1\\
+167{,}6
\end{array}
\]

Diese Differenzen sind noch kleiner als die mit dem Integral, sie scheinen aber bei zunehmendem $n$ schneller zu wachsen als diese, so dass leicht möglich
\SourceLeaf{456}%
wäre, dass bei viel weiterer Fortsetzung jene die letztern überträfen. Um Zählung und Formel in Uebereinstimmung zu bringen, müsste man respective anstatt $A=1{,}08366$ setzen
\[
\begin{array}{r}
1{,}09040\\
1{,}07682\\
1{,}07582\\
1{,}07529\\
1{,}07179\\
1{,}07297
\end{array}
\]

Es scheint, dass bei wachsendem $n$ der (Durchschnitts-)Werth von $A$ abnimmt, ob aber die Grenze beim Wachsen des $n$ ins Unendliche $1$ oder eine von $1$ verschiedene Grösse sein wird, darüber wage ich keine Vermuthung. Ich kann nicht sagen, dass eine Befugniss da ist, einen ganz einfachen Grenzwerth zu erwarten; von der andern Seite könnte der Ueberschuss des $A$ über $1$ ganz füglich eine Grösse von der Ordnung $\frac{1}{\log n}$ sein. Ich würde geneigt sein zu glauben, dass das Differential der betreffenden Function einfacher sein muss, als die Function selbst. Indem ich für jene $\frac{\mathrm{d}n}{\log n}$ vorausgesetzt habe, würde \textsc{Legendre}’s Formel eine Differentialfunction voraussetzen, die etwa $\frac{\mathrm{d}n}{\log n-(A-1)}$ wäre. Ihre Formel übrigens würde für ein sehr grosses $n$ als mit
\[
\frac{n}{\log n-\frac{1}{2k}}
\]
übereinstimmend betrachtet werden können, wo $k$ der Modulus der \textsc{Brig\-gi}’schen Logarithmen ist, also mit \textsc{Legendre}’s Formel, wenn man
\[
A=\tfrac{1}{2k}=1{,}1513\ \text{setzt.}
\]

Endlich will ich noch bemerken, dass ich zwischen Ihren Abzählungen und den meinigen ein Paar Differenzen bemerkt habe.
\begin{center}
\begin{tabular}{lrlrlrlr}
Zwischen & 59000 & u. & 60000 & haben Sie & 95 & ich & 94\\
 & 101000 & & 102000 & & 94 & & 93
\end{tabular}
\end{center}
Die erste Differenz hat vielleicht ihren Grund darin, dass in \textsc{Lambert}’s Suppl. die Primzahl 59023 zweimal aufgeführt ist. Die Chiliade von 101000—102000 wimmelt in \textsc{Lambert}’s Supplementen von Fehlern, ich habe in meinem Exemplare 7 Zahlen angestrichen, die keine Primzahlen sind, und dagegen 2 fehlende ein%
\SourceLeaf{457}%
geschaltet. Könnten Sie nicht den jungen \textsc{Dase} veranlassen, dass er die Primzahlen in den folgenden Millionen aus denjenigen bei der Akademie befindlichen Tafeln abzählte, die wie ich fürchte das Publicum nicht besitzen soll? Für diesen Fall bemerke ich, dass in der 2. und 3. Million die Abzählung auf meine Vorschrift nach einem besondern Schema gemacht ist, welches ich selbst auch schon bei einem Theile der ersten Million angewandt hatte. Die Abzählungen von je 100000 stehen auf Einer (klein) Octavseite in 10 Columnen, jede sich auf Eine Myriade beziehend; dazu kommt noch eine Columne davor (links) und eine dahinter rechts; als Beispiel hier eine Verticalcolumne und die beiden Zusatzcolumnen aus dem Intervall $1000000\mathinner{.\,.\,.}1100000$ \mbox{— — —}

Zur Erläuterung diene z. B. die 1. Verticalreihe. In der Myriade 1000000 bis 1010000 sind 100 Hecatontaden; darunter ist 1 die nur eine Primzahl enthält; gar keine mit 2 oder 3; 2 Stück mit je 4 Primzahlen; 11 Stück mit je 5 u. s. w. alle zusammen geben $752=1\pdot1+4\pdot2+5\pdot11+6\pdot14+\dotstwo$ Die letzte Columne enthält die Aggregate aus den 10 einzelnen. Die Zahlen 14. 15. 16 in der ersten Verticalreihe stehen hier nur zum Ueberfluss, da keine Hecatontaden mit so vielen Primzahlen vorkommen; aber auf den folgenden Blättern bekommen sie Geltung. Zuletzt werden wieder die 10 Seiten in 1 vereinigt, und umfassen so die ganze 2te Million.

Doch es ist Zeit abzubrechen. \mbox{— — —} Unter herzlichen Wünschen für Ihr Wohlbefinden
\begin{center}Stets der Ihrige\end{center}
\noindent Göttingen, 24. December 1849.\hfill C. F. \textsc{Gauss}.
\BookRule{.25\textwidth}
\par\endgroup
